% GENERATED FILE. Edit canonical lessons, then run npm run generate:books.
% canonical-source-hash: fnv1a64:e94150cebf39ad6c
% canonical-lessons: SA-C175-harati, SA-C175-bharati, SA-C175-vahati, SA-C175-sahate, SA-C175-yatate

\chapter{Takes away, Supports, Carries, Endures, Tries}
\label{ch:sa-takes-away-supports-carries-endures-tries}
\hlchaptermodality{175}

\begin{chapteropening}
\textbf{By the end of this chapter:} \emph{I can say takes away, supports, carries, endures and tries.}

Complete the last lesson of chapter 175: I can say takes away, supports, carries, endures and tries.
\end{chapteropening}

\section[\texorpdfstring{\sk{हरति} (harati)}{हरति (harati)}]{\sk{हरति} (harati) — takes away, carries off}

\label{lesson:SA-C175-harati}



\begin{quote}
Before the new one: say the Sanskrit for keeps, then the Sanskrit for hurts.
\end{quote}

\subsection*{You'll want to know: \sk{हरति}}
\textbf{\sk{हरति}} — \emph{harati} — \textquotedblleft{}takes away, carries off\textquotedblright{}.

\begin{grammarlens}[title={What you've built}]
One more word for this chapter.
\end{grammarlens}

\subsection*{Your turn}
\begin{itemize}
\raggedright
  \item \emph{Say it:} \emph{harati}
  \item \emph{Say it:} \emph{harati}, once more
  \item \emph{Say it:} say \emph{harati}
  \item \emph{Recall:} say the Sanskrit for keeps, then the Sanskrit for hurts, then say \emph{harati} again
\end{itemize}

\begin{tcolorbox}[breakable,colback=teal!4,colframe=teal!35!black,title={Before you move on}]
What does \sk{हरति} mean? (\textquotedblleft{}Takes away, carries off\textquotedblright{}.) Say it once more.
\end{tcolorbox}
\section[\texorpdfstring{\sk{भरति} (bharati)}{भरति (bharati)}]{\sk{भरति} (bharati) — supports, bears}

\label{lesson:SA-C175-bharati}



\begin{quote}
Before the new one: say the Sanskrit for hurts, then the Sanskrit for takes away.
\end{quote}

\subsection*{You'll want to know: \sk{भरति}}
\textbf{\sk{भरति}} — \emph{bharati} — \textquotedblleft{}supports, bears\textquotedblright{}.

\begin{grammarlens}[title={What you've built}]
One more word for this chapter.
\end{grammarlens}

\subsection*{Your turn}
\begin{itemize}
\raggedright
  \item \emph{Say it:} \emph{bharati}
  \item \emph{Say it:} \emph{bharati}, once more
  \item \emph{Say it:} say \emph{bharati}
  \item \emph{Recall:} say the Sanskrit for hurts, then the Sanskrit for takes away, then say \emph{bharati} again
\end{itemize}

\begin{tcolorbox}[breakable,colback=teal!4,colframe=teal!35!black,title={Before you move on}]
What does \sk{भरति} mean? (\textquotedblleft{}Supports, bears\textquotedblright{}.) Say it once more.
\end{tcolorbox}
\section[\texorpdfstring{\sk{वहति} (vahati)}{वहति (vahati)}]{\sk{वहति} (vahati) — carries, flows}

\label{lesson:SA-C175-vahati}



\begin{quote}
Before the new one: say the Sanskrit for takes away, then the Sanskrit for supports.
\end{quote}

\subsection*{You'll want to know: \sk{वहति}}
\textbf{\sk{वहति}} — \emph{vahati} — \textquotedblleft{}carries, flows\textquotedblright{}.

\begin{grammarlens}[title={What you've built}]
One more word for this chapter.
\end{grammarlens}

\subsection*{Your turn}
\begin{itemize}
\raggedright
  \item \emph{Say it:} \emph{vahati}
  \item \emph{Say it:} \emph{vahati}, once more
  \item \emph{Say it:} say \emph{vahati}
  \item \emph{Recall:} say the Sanskrit for takes away, then the Sanskrit for supports, then say \emph{vahati} again
\end{itemize}

\begin{tcolorbox}[breakable,colback=teal!4,colframe=teal!35!black,title={Before you move on}]
What does \sk{वहति} mean? (\textquotedblleft{}Carries, flows\textquotedblright{}.) Say it once more.
\end{tcolorbox}
\section[\texorpdfstring{\sk{सहते} (sahate)}{सहते (sahate)}]{\sk{सहते} (sahate) — endures, bears}

\label{lesson:SA-C175-sahate}



\begin{quote}
Before the new one: say the Sanskrit for supports, then the Sanskrit for carries.
\end{quote}

\subsection*{You'll want to know: \sk{सहते}}
\textbf{\sk{सहते}} — \emph{sahate} — \textquotedblleft{}endures, bears\textquotedblright{}.

\begin{grammarlens}[title={What you've built}]
One more word for this chapter.
\end{grammarlens}

\subsection*{Your turn}
\begin{itemize}
\raggedright
  \item \emph{Say it:} \emph{sahate}
  \item \emph{Say it:} \emph{sahate}, once more
  \item \emph{Say it:} say \emph{sahate}
  \item \emph{Recall:} say the Sanskrit for supports, then the Sanskrit for carries, then say \emph{sahate} again
\end{itemize}

\begin{tcolorbox}[breakable,colback=teal!4,colframe=teal!35!black,title={Before you move on}]
What does \sk{सहते} mean? (\textquotedblleft{}Endures, bears\textquotedblright{}.) Say it once more.
\end{tcolorbox}
\section[\texorpdfstring{\sk{यतते} (yatate)}{यतते (yatate)}]{\sk{यतते} (yatate) — tries, strives}

\label{lesson:SA-C175-yatate}



\begin{quote}
Before the new one: say the Sanskrit for carries, then the Sanskrit for endures.
\end{quote}

\subsection*{You'll want to know: \sk{यतते}}
\textbf{\sk{यतते}} — \emph{yatate} — \textquotedblleft{}tries, strives\textquotedblright{}.

\begin{grammarlens}[title={What you've built}]
One more word for this chapter.
\end{grammarlens}

\subsection*{Your turn}
\begin{itemize}
\raggedright
  \item \emph{Say it:} \emph{yatate}
  \item \emph{Say it:} \emph{yatate}, once more
  \item \emph{Say it:} say \emph{yatate}
  \item \emph{Recall:} say the Sanskrit for carries, then the Sanskrit for endures, then say \emph{yatate} again
\end{itemize}

\begin{tcolorbox}[breakable,colback=teal!4,colframe=teal!35!black,title={Before you move on}]
What does \sk{यतते} mean? (\textquotedblleft{}Tries, strives\textquotedblright{}.) Say it once more.
\end{tcolorbox}
